\documentclass[10pt]{exam}
\usepackage[hon]{template-for-exam}
\usepackage{fancybox}

\title{Newton's Laws of Motion: Reading Guide}
\author{Rohrbach}
\date{\today}

\begin{document}
\maketitle

% \section*{Definitions}

% \begin{itemize}
%   \item kinematics
%   \item directly proportional
%   \item inversely proportional
%   \item weight
% \end{itemize}

% \pagebreak


\section*{4.1 Development of Force Concept}

\subsection*{Definitions}

\begin{itemize}
  \item dynamics
  \item force
  \item free-body diagram
\end{itemize}

\subsection*{Concept Check}

A force is which of these?

\begin{oneparcheckboxes}
  \correctchoice a vector
  \choice a scalar
  \choice a quasar
  \choice a quaternion
\end{oneparcheckboxes}

\vspace{3em}
\hrule

\section*{4.2 Newton's First Law of Motion: Inertia}

\begin{Sbox}
  \begin{minipage}{6in}
    \subsection*{Newton's First Law (NFL)}

    A body at rest \fillin[remains][7em] at rest, or, if in motion, \fillin[remains][5em] at a \fillin[constant velocity][10em] unless acted on by a \fillin[net external force][15em].
  \end{minipage}
\end{Sbox}

\shadowbox{\TheSbox}

\vspace{1em}

\subsection*{Thinking about the First Law}

According to the book, what is the important verb in Newton's First Law, and why is it important?

\vspace{5em}


\vspace{1em}

\subsection*{Concept Check}

\paragraph{A trick question!} A school bus comes to a sudden stop, and all of the backpacks on the floor start to slide forward.  What force causes them to do that?

\begin{oneparcheckboxes}
  \choice the force of the brakes
  \choice friction
  \choice the force of inertia
  \correctchoice there is no force that causes that!
\end{oneparcheckboxes}

\hfill

Explain your answer:

\vspace{3em}

\subsection*{Definitions}

\begin{itemize}
  \item inertia
  \item mass
\end{itemize}



\section*{4.3 Newton's Second Law of Motion}

\begin{Sbox}
  \begin{minipage}{6in}
    \subsection*{Newton's Second Law (NSL)}

    The acceleration of an object is
    %
    \begin{itemize}
      \item directly proportional to and in the same \fillin[direction][10em] as the \fillin[the net external force][10em] acting on it.
      \item inversely proportional to its \fillin[mass][5em]
    \end{itemize}
    \paragraph{Equations:}\hfill
    \vspace{3em}
  \end{minipage}
\end{Sbox}

\shadowbox{\TheSbox}

\subsection*{Definitions}

\begin{itemize}
  \item inversely proportional
  \item directly proportional
  \item weight
\end{itemize}

\subsection*{Weight and Mass}
Summarize what weight is and how it is different than mass.

\vspace{6em}

Equation for weight:

\vspace{2em}


\subsection*{Concept Check}

An astronaut on the moon has...

\begin{oneparcheckboxes}
  \choice less mass than on earth
  \correctchoice less weight than on earth
  \choice less mass and less weight than on earth
\end{oneparcheckboxes}

\hfill

Explain your answer:

\vspace{3em}

\vspace{1em}

\hrule

\pagebreak

\section*{4.4 Newton's Third Law of Motion: Symmetry in Forces}

\begin{Sbox}
  \begin{minipage}{6in}
    \subsection*{Newton's Third Law (NTL)}

    Whenever one body exerts a force on a second body, the first body experiences a force that is \fillin[equal][7em] in magnitude and \fillin[opposite][7em] in direction to the force that it exerts.

    \paragraph{Equation:}\hfill 

    \vspace{3em}

    \paragraph{Shortened form:} Every \fillin[action][7em] has an \fillin[equal][6em] and \fillin[opposite][6em] reaction.

  \end{minipage}
\end{Sbox}

\shadowbox{\TheSbox}

\subsection*{Examples}
Write down three examples of Newton's Third Law in action from the book.

\vs

\subsection*{Concept Check}

A hammer hits a nail. If the action force is the hammer pushes the nail down, what is the reaction force?

\hfill

\begin{checkboxes}
  \choice The nail moves down
  \choice The wood pushes the nail up
  \choice The hammer pushes the nail up
  \choice The gravity of the hammer pulls it down
  \correctchoice The nail pushes the hammer up
\end{checkboxes}

\hfill

\hfill

\noindent
An airplane experiences a force of gravity in which the earth pulls the plane down.  What is the reaction force to this?

\hfill

\begin{checkboxes}
  \choice The airplane's wings push the plane up
  \correctchoice The airplane pulls the earth up
  \choice The airplane's wings push the air down
  \choice The airplane eventually comes down.
\end{checkboxes}

\pagebreak \hfill









\end{document}